\section{Turing Machine --- NDCodes}
%%% generated by the blueprint pipeline from TCSlib.Complexity.TuringMachine.NDCodes
\subsection*{Turing Machine --- NDCodes: Introduction}

Declarations of the Lean module \texttt{TCSlib.Complexity.TuringMachine.NDCodes}.

\subsection*{Turing Machine --- NDCodes: Results}

\begin{definition}[Two-tape nondeterministic machine code]
\lean{Turing.CodeNDTM}
\label{Turing.CodeNDTM}
\leanok
\uses{Turing.NDTM}
A code for a nondeterministic Turing machine consists of a natural number $n$ and a
nondeterministic machine with two work tapes, binary tape symbols, and state set
$\{0,\ldots,n\}$; thus the state-count parameter is one less than the number of states.
\end{definition}

\begin{definition}[Finite machine associated to a coded machine]
\lean{Turing.CodeNDTM.toFinNDTM}
\label{Turing.CodeNDTM.toFinNDTM}
\leanok
\uses{Turing.CodeNDTM, Turing.FinNDTM}
For a coded nondeterministic machine with two work tapes, associate the finite machine
with the same transition behavior, two work tapes, and state set consisting of the
indices $0,\ldots,n$, where $n$ is the coded machine’s state parameter.
\end{definition}

\begin{definition}[Binary serialization of a coded machine]
\lean{Turing.CodeNDTM.serialize}
\label{Turing.CodeNDTM.serialize}
\leanok
\uses{Turing.CodeNDTM, Turing.actionBits₂, Turing.pairEncode, Turing.unaryFin, Turing.workPair}
For a coded nondeterministic machine with $n+1$ states, serialize its state parameter in
binary and pair that string with a bit string containing its initial state in unary,
followed by the transition actions. The actions are listed for both choice bits, every
state, and every combination of input and two work-tape symbols, with each symbol either
blank, $0$, or $1$.
\end{definition}

\begin{definition}[Encoding and decoding machine codes]
\lean{Turing.NDMachineCode}
\label{Turing.NDMachineCode}
\leanok
\uses{Turing.CodeNDTM}
A nondeterministic machine coding scheme assigns a binary string to each coded machine
and assigns a coded machine to every binary string. For every coded machine $M$ and
every natural number $m$, decoding the encoding of $M$ followed by $m$ copies of the bit
$1$ returns $M$.
\end{definition}

\begin{theorem}[Decoding an encoded machine]
\lean{Turing.NDMachineCode.decode_encode}
\label{Turing.NDMachineCode.decode_encode}
\leanok
\uses{Turing.CodeNDTM, Turing.NDMachineCode}
\difficulty{1}
For every machine code scheme $c$ satisfying the property that decoding the encoding of
any machine followed by any number of padding bits equal to $1$ returns that machine,
and every encoded machine $M$, decoding the encoding of $M$ returns $M$.
\end{theorem}

\begin{definition}[Effective nondeterministic machine coding]
\lean{Turing.EffectiveNDMachineCode}
\label{Turing.EffectiveNDMachineCode}
\leanok
\uses{Turing.CodeNDTM.serialize, Turing.FinTM, Turing.FinTM.ComputesFunInTime, Turing.NDMachineCode}
An effective nondeterministic machine coding scheme consists of a machine coding scheme,
a deterministic Turing machine that computes a function within a specified time bound,
and a time bound assigning a natural number to each input length. The computed function
sends each binary string to the fixed serialization of the machine obtained by decoding
that string.
\end{definition}

\begin{definition}[Binary format for coded nondeterministic machines]
\lean{Turing.ndFormat}
\label{Turing.ndFormat}
\leanok
\uses{Turing.CodeFormat, Turing.CodeNDTM, Turing.CodeNDTM.serialize, Turing.CodeTable2, Turing.actionBits₂_eq_actionBitsK, Turing.codeEraseTable2, Turing.codeErase_bind, Turing.codeReadTable2, Turing.codeReadTable2_append, Turing.codeReadTable2_sound, Turing.codeSkipRepeat, Turing.codeSkipTable2, Turing.codeTable2Bits, Turing.codeTable2Bits_length, Turing.workPair}
A code format for two-tape nondeterministic machines uses, for each state count, a pair
of transition tables, one for each choice bit. It serializes a machine by pairing the
binary expansion of its state parameter with the unary initial-state index followed by
the two table encodings; its table reader recovers a table and the remaining bits, and
its erasing reader discards the table encoding. The format has guard constant $702$ and
uses a fixed fallback machine for malformed strings.
\end{definition}

\begin{theorem}[Existence of a nondeterministic effective code scheme]
\lean{Turing.exists_effectiveNDMachineCode}
\label{Turing.exists_effectiveNDMachineCode}
\leanok
\uses{Turing.CodeFormat.decode, Turing.CodeFormat.decode_serialize_pad, Turing.CodeFormat.exists_canonizer, Turing.CodeNDTM.serialize, Turing.EffectiveNDMachineCode, Turing.codePrimRepeat, Turing.codePrimSkipTable2, Turing.codeSkipTable2, Turing.ndFormat}
\difficulty{4}
There exists an effective coding scheme for nondeterministic two-work-tape Turing
machines over the binary alphabet: it has an encoding and a total decoder such that, for
every coded machine $M$ and every natural number $m$, decoding the encoding of $M$
followed by $m$ copies of the bit $1$ returns $M$. The scheme also includes a finite
binary-alphabet Turing machine and a time bound that, on every binary input string,
compute the fixed serialization of the machine decoded from that string.
\end{theorem}
